%  LaTeX support: latex@mdpi.com 
%  For support, please attach all files needed for compiling as well as the log file, and specify your operating system, LaTeX version, and LaTeX editor.

%=================================================================
\documentclass[land,article,submit,pdftex,moreauthors]{Definitions/mdpi} 



% MDPI internal commands
\firstpage{1} 
\makeatletter 
\setcounter{page}{\@firstpage} 
\makeatother
\pubvolume{1}
\issuenum{1}
\articlenumber{0}
\pubyear{2023}
\copyrightyear{2023}
%\externaleditor{Academic Editor: Firstname Lastname}
\datereceived{2 July 2023} 
\daterevised{} 
\dateaccepted{} 
\datepublished{} 
\hreflink{https://doi.org/} % If needed use \linebreak
%\doinum{}

%=================================================================
% Add packages and commands here. The following packages are loaded in our class file: fontenc, inputenc, calc, indentfirst, fancyhdr, graphicx, epstopdf, lastpage, ifthen, lineno, float, amsmath, setspace, enumitem, mathpazo, booktabs, titlesec, etoolbox, tabto, xcolor, soul, multirow, microtype, tikz, totcount, changepage, attrib, upgreek, cleveref, amsthm, hyphenat, natbib, hyperref, footmisc, url, geometry, newfloat, caption

\graphicspath{{figures/}} % path to folder with figures
\usepackage{subcaption}
\usepackage{listings}
\usepackage{xcolor}
\usepackage{gensymb} % for \textdegree
\usepackage{tabularx}

%=================================================================
% Full title of the paper (Capitalized)
\Title{Seafloor and Ocean Crust Structure of the Kerguelen Plateau from Marine Geophysical and Satellite Altimetry Datasets}
% Interplay Between Seafloor Spreading, Marine Gravity Roughness and Ocean Crust Structure of the Kerguelen Plateau Analysed from Satellite Altimetry

% MDPI internal command: Title for citation in the left column
\TitleCitation{Seafloor and Ocean Crust Structure of the Kerguelen Plateau from Marine Geophysical and Satellite Altimetry Datasets}

% Author Orchid ID: enter ID or remove command
\newcommand{\orcidauthorA}{0000-0002-5759-1089} % Polina Add \orcidA{} behind the author's name

% Authors, for the paper (add full first names)
\Author{Polina Lemenkova $^{1}$*\orcidA{}}

%\longauthorlist{yes}

% MDPI internal command: Authors, for metadata in PDF
\AuthorNames{Polina Lemenkova}

% MDPI internal command: Authors, for citation in the left column
\AuthorCitation{Lemenkova, P.}
% If this is a Chicago style journal: Lastname, Firstname, Firstname Lastname, and Firstname Lastname.

% Affiliations / Addresses (Add [1] after \address if there is only one affiliation.)
\address{%
$^{1}$ \quad Laboratory of Image Synthesis and Analysis (LISA), École polytechnique de Bruxelles (Brussels Faculty of Engineering), Université Libre de Bruxelles (ULB). Building L, Campus du Solbosch, ULB -- LISA CP165/57, Avenue Franklin D. Roosevelt 50, Brussels 1050, Belgium.
}

% Contact information of the corresponding author
\corres{Correspondence: polina.lemenkova@ulb.be; Tel.: +32 471 86 04 59 (P.L.)}

% Current address and/or shared authorship
%\firstnote{Current address: Affiliation 3.} 

% Abstract (Do not insert blank lines, i.e. \\) 
\abstract{The volcanic Kerguelen Islands are formed on one of the world's largest submarine plateaus. Located in the remote segment of the southern Indian Ocean close to Antarctica, the Kerguelen Plateau is notable for complex tectonic origin and geologic formation related to the Cretaceous history of the continents. This is reflected in varying age of the oceanic crust adjacent to the plateau and highly heterogeneous bathymetry of the Keguelen Plateau where seafloor structure of the southern segment  differs significantly from its northern part. Remote sensing data derived from marine gravity and satellite radar altimetry surveys serve as an important source of information for mapping of such complex seafloor features. Automatic matching of topographic and geophysical grids with high accuracy is essential for complex geologic-tectonic investigations. We propose incorporating the geospatial information from NOAA, EMAG2 and WDMAM, ETOPO1, EGM96 and EGM2008 to refine the extent and distribution of features extracted from geospatial datasets. The cartographic joint analysis of topography, magnetic anomalies, tectonic and gravity data is based on integrated mapping using scripting toolset Generic Mapping Tools (GMT). The refined mapping of the submerged features (Broken Ridge orientation, seamounts around the Crozet Islands, fabric of seafloor and orientation and frequency of magnetic anomalies) is used to analyse the correspondence of the objects with regard to free-air gravity and magnetic anomalies, geodynamics and seabed structure in south-west Indian Ocean. Experimental results show that integrating the proposed datasets using the state-of-the-art cartographic scripting language improves identification and visualization of the seabed objects. The presented research includes 10 new maps which contribute to  increasing the knowledge of subsurface seafloor structure around the Kerguelen Plateau and composition of the Earth's crust in the French Southern and Antarctic Lands.
}

% Keywords
\keyword{Antarctic; Southern Ocean; bathymetry; French Southern and Antarctic Lands; cartography; satellite altimetry; marine geophysics; keyword 8; keyword 9; keyword 10} 

% The fields PACS, MSC, and JEL may be left empty or commented out if not applicable
%\PACS{J0101}
%\MSC{}
%\JEL{}

\begin{document}

%----------- section ----------->
\section{Introduction}

Analysis of seafloor structure and ocean crust concerns mapping, detecting, and recognising the distribution of bathymetric structures and variations in geophysical fields analysed from remote sensing data. Among these problems, integrated geophysics interpretation of satellite altimetry,  marine gravimetry and magnetic anomalies is one of the most active research areas and has attracted much research interest over the past decades. Previous studies \cite{ESHAGH2021313} provided a broad overview of numerous approaches for analysing satellite gravimetry data in marine geophysical research. Marine geophysical and satellite altimetry data provide information on the density in the Earth's crust and processes in the upper mantle \cite{KAPAWAR2021106692}, sediment thickness and basement depths \cite{Lemenkova202005}. Moreover,  the analysis of the geophysical data enables to reveal the key characteristics regarding the lithosphere thickness such as Moho discontinuity, elastic thickness \cite{TESAURO201378}, viscosity and flexural rigidity. Such data are essential for our understanding of the Earth's structure.

In terms of the data-driven analyses, previous works investigating the oceanic seabed can be roughly classified into studies focused on bathymetric mapping and geophysical analysis. Mapping seafloor usually characterises the topographic patterns of the oceanic crust with increased cartographic details of data visualization and improved methods \cite{Battista,Lemenkova202105,Lemenkova202012}. Marine geophysical methods generally investigate the geology of continental margins with deep-sea drilling \cite{Schlich1992TheGA}, analysis of features extracted from seismic survey and remote sensing data \cite{NEGI2022229330}, investigating the structure of the deep ocean basins \cite{BENARD2010633,RAMANA2001241,DAVIS201616} and modelling mid-ocean ridge system in a global or regional context \cite{FREY200073,Sreejith,Lemenkova202112}.

\pagebreak
\begin{figure}[H]
\makebox[1.0\linewidth][r]{%
\centering
	\hspace{100pt}\includegraphics[width=14.0 cm]{Fig_01.jpg}
}
\caption{Map of the Kerguelen Plateau region, Southern Ocean. Mapping software: GMT v. 6-1-1. Map source: author. \label{fig01}}
\end{figure}

%----------- section ----------->
\section{Materials and Methods}

%----------- subsection ----------->
\subsection{Subsection}

%----------- subsection ----------->
\subsection{Subsection}

%----------- subsection ----------->
\subsection{Subsection}

%----------- section ----------->
\pagebreak
\section{Results}

%----------- subsection ----------->
\subsection{Ocean floor formation}

Age, spreading rates and spreading symmetry of the ocean crust indicate the evolution steps in ocean floor formation in South-West Indian Ocean and the Kerguelen Plateau. The age of the ocean crust was determined by interpolation of adjacent seafloor isochrons oriented towards the direction of seafloor spreading, Figure \ref{fig02}.

\begin{figure}[H]
\makebox[1.0\linewidth][r]{%
\centering
	\hspace{100pt}\includegraphics[width=14.0 cm]{Fig_02.jpg}
}
\caption{Age of the ocean crust in SW Indian Ocean and the Kerguelen Plateau region and east Antarctic. Mapping software: GMT v. 6-1-1. Map source: author. \label{fig02}}
\end{figure}


\begin{figure}[H]
\makebox[1.0\linewidth][r]{%
\centering
	\hspace{100pt}\includegraphics[width=14.0 cm]{Fig_03.jpg}
}
\caption{Age, spreading rates and spreading symmetry of the ocean crust over the Kerguelen Plateau region, east Antarctic and south-west Indian Ocean: asymmetries in crustal accretion (\%) on conjugate ridge flanks. Mapping software: GMT v. 6-1-1. Map source: author. \label{fig03}}
\end{figure}

Spreading half rates of the ocean crust are visualised based on the 2 arc min netCDF grid v.3.6 (NOAA). 

\begin{figure}[H]
\makebox[1.0\linewidth][r]{%
\centering
	\hspace{100pt}\includegraphics[width=14.0 cm]{Fig_04.jpg}
}
\caption{Spreading half rates of the ocean crust over the Kerguelen Plateau region, east Antarctic and south-west Indian Ocean with added GSFML data. Mapping software: GMT v. 6-1-1. Map source: author. \label{fig04}}
\end{figure}

\begin{figure}[H]
\makebox[1.0\linewidth][r]{%
\centering
	\hspace{100pt}\includegraphics[width=14.0 cm]{Fig_05.jpg}
}
\caption{Age uncertainty in lithosphere crust (m.y) over the Kerguelen Plateau region, east Antarctic and south-west Indian Ocean. Cartography: GMT, version 6-1-1. Map source: author. \label{fig05}}
\end{figure}

%----------- subsection ----------->
\subsection{Sediment thickness}

\begin{figure}[H]
\makebox[1.0\linewidth][r]{%
\centering
	\hspace{100pt}\includegraphics[width=14.0 cm]{Fig_06.jpg}
}
\caption{Sediment thickness (m) in East Antarctica, South-West Indian Ocean and the Kerguelen Plateau region. Cartography: GMT. Lambert azimuthal equal-area projection. Map source: author. \label{fig06}}
\end{figure}

%----------- subsection ----------->
\subsection{Free-air gravity anomaly and vertical gravity gradients}

\begin{figure}[H]
\makebox[1.0\linewidth][r]{%
\centering
	\hspace{100pt}\includegraphics[width=14.0 cm]{Fig_07.jpg}
}
\caption{Marine gravity field over the Kerguelen Plateau region, south-west Indian Ocean. Mapping software: GMT v. 6-1-1. Map source: author. \label{fig07}}
\end{figure}

\begin{figure}[H]
\makebox[1.0\linewidth][r]{%
\centering
	\hspace{100pt}\includegraphics[width=14.0 cm]{Fig_08.jpg}
}
\caption{Vertical gravity gradient anomalies over the Kerguelen Plateau region showing variations in gravity with the change in topographic elevations. Mapping software: GMT v. 6-1-1. Map source: author. \label{fig08}}
\end{figure}

%----------- subsection ----------->
\subsection{Magnetic anomalies over Kerguelen Plateau}
The anomaly of the magnetic intensity at an altitude of 5 km above mean sea level over Kerguelen is shown in Figure \ref{fig09}

\begin{figure}[H]
\makebox[1.0\linewidth][r]{%
	\begin{subfigure}[b]{.65\textwidth}
		\centering
			\includegraphics[width=9.0cm]{Fig_09a.jpg}
		\caption{}
	\end{subfigure}	
	%\vspace{1mm}
	\begin{subfigure}[b]{.60\textwidth}
		\centering
			\includegraphics[width=9.5cm]{Fig_09b.jpg}
		\caption{}
	\end{subfigure}%
}
\caption{(\textbf{a}): Marine and Earth airborne magnetic anomaly grid (3-min resolution) based on WDMAM, south-west Indian Ocean; (\textbf{b}): Marine and Earth airborne magnetic anomaly grid based on EMAG-2 over the Kerguelen Plateau region. Black areas signify "no data" in the original grid. Cartography: GMT, version 6-1-1. Map source: author.}\label{fig09}
\end{figure}

%----------- subsection ----------->
\subsection{Geoid models}
The variations in the geopotential geoid models over Kerguelen is shown in Figure \ref{fig10}

\begin{figure}[H]
\makebox[1.0\linewidth][r]{%
	\begin{subfigure}[b]{.75\textwidth}
		\centering
			\includegraphics[width=9.6cm]{Fig_10a.jpg}
		\caption{}
	\end{subfigure}	\\
	
	\begin{subfigure}[b]{.60\textwidth}
		\centering
			\includegraphics[width=9.5cm]{Fig_10b.jpg}
		\caption{}
	\end{subfigure}%
}
\caption{(\textbf{a}): Geoid model based on EGM-96 over the Kerguelen Plateau region, east Antarctic and south-west Indian Ocean; (\textbf{b}): Geoid model based on EGM-2008 over the Kerguelen Plateau region, east Antarctic and south-west Indian Ocean. Cartography: GMT, version 6-1-1. Map source: author.}\label{fig10}
\end{figure}

%----------- section ----------->
\section{Discussion}

%----------- section ----------->
\section{Conclusions}

%----------- section ----------->
\authorcontributions{Supervision, conceptualization, methodology, software, resources, funding acquisition, and project administration, writing---original draft preparation, methodology, software, data curation, visualization, formal analysis, validation, writing---review and editing, and investigation, P.L.}

\funding{The publication was funded by the Editorial Office of \emph{Land}, Multidisciplinary Digital Publishing Institute (MDPI), by providing 100\% discount for the APC of this manuscript.}

\institutionalreview{Not applicable.}

\informedconsent{Not applicable.}

\dataavailability{Not applicable} 

\acknowledgments{The authors thank the reviewers for reading and review of this manuscript.}

\conflictsofinterest{The authors declare no conflict of interest.} 

\abbreviations{Abbreviations}{
The following abbreviations are used in this manuscript:\\

\noindent 
\begin{tabular}{@{}ll}
MDPI & Multidisciplinary Digital Publishing Institute\\
DOAJ & Directory of open access journals\\
LD & Linear dichroism
\end{tabular}
}

%%%%%%%%%%%%%%%%%%%%%%%%%%%%%%%%%%%%%%%%%%
%% Optional
\appendixtitles{no} % Leave argument "no" if all appendix headings stay EMPTY (then no dot is printed after "Appendix A"). If the appendix sections contain a heading then change the argument to "yes".
\appendixstart
\appendix
\section[\appendixname~\thesection]{}
\subsection[\appendixname~\thesubsection]{}
The appendix 

\section[\appendixname~\thesection]{}
All appendix sections must be cited in the main text. In the appendices, Figures, Tables, etc. should be labeled, starting with ``A''---e.g., Figure A1, Figure A2, etc.

%%%%%%%%%%%%%%%%%%%%%%%%%%%%%%%%%%%%%%%%%%
\begin{adjustwidth}{-\extralength}{0cm}
%\printendnotes[custom] % Un-comment to print a list of endnotes

\reftitle{References}
%\nocite{*}

%=====================================
% References, variant A: external bibliography
%=====================================
\bibliography{Kerguelen.bib}

%%%%%%%%%%%%%%%%%%%%%%%%%%%%%%%%%%%%%%%%%%
\end{adjustwidth}
\end{document}
